\subsection{不可约表示}

本节中，G 是拓扑群。

定义 2. --- 拓扑 K-向量空间 E 中 G 的表示 \(\varrho\) 称为不可约的，是指 \(\{0\}\) 在 E 中不稠密，并且 \(\varrho\) 的唯一子表示是 \(\varrho\) 本身和在 \(\{0\}\) 的闭包中诱导的表示。

如果 \(\varrho\) 是分离拓扑向量空间 E 中 G 的不可约表示，那么 E 的每个非零元素都是 \(\varrho\) 的循环向量。

引理 2. --- 设 \(\pi\) 和 \(\varrho\) 是分离拓扑 K-向量空间中 G 的连续线性表示。假设 \(\pi\) 不可约。从 \(\pi\) 到 \(\varrho\) 的每个非零 G-态射都是单射，而从 \(\varrho\) 到 \(\pi\) 的每个非零 G-态射都有稠密像。

特别地，如果 \(\pi\) 和 \(\varrho\) 都不可约且有限维，那么从 \(\pi\) 到 \(\varrho\) 的每个非零 G-态射都是同构。

由于 \(\varrho\) 的表示空间是分离的，G-态射 \(u\)（从 \(\pi\) 到 \(\varrho\)）的核是闭的，因而定义 \(\pi\) 的一个子表示。如果态射 \(u\) 非零，其核必为 \(\{0\}\)，因为 \(\pi\) 是分离拓扑向量空间中的不可约表示。同样，从 \(\varrho\) 到 \(\pi\) 的非零 G-态射的像的闭包是 \(\pi\) 的一个非零子表示，因而等于 \(\pi\) 的表示空间。

最后一个断言由上述内容推出。

引理 3. --- 设 \(\varrho\) 是分离拓扑 K-向量空间 E 中的有限维连续线性表示。如果 E 非零，那么 \(\varrho\) 存在一个不可约子表示。

由于 E 有限维且非零，存在维数最小的非零 G-不变子空间 F。该子空间在 E 中闭，并定义 E 的一个子表示；F 的每个子表示也是 E 的子表示，因此由极小性，F 中的表示不可约。

注. --- G 的一个非零表示 E 并不总是包含不可约子表示（参见 V, p. 426, 注）。不过将会看到，当 G 是紧群、K = C，且 E 是 K 上非零的分离拟完备局部凸空间时，情况确实如此（V, p. 464 的命题 7）。

